\section{Ontology Design and Schema Modeling}
\label{sec:gr-ontology-design}

An ontology supplies the schema that the extracted facts must conform to,
specifying the conceptualization of the domain in terms of classes, properties,
and constraints~\cite{gruber1993ontology}. Committing to such a shared vocabulary
is what gives the knowledge graph consistent meaning and enables validation and
inference over it~\cite{gruber1993ontology, hogan2021kg}. In a Graph RAG system
the ontology therefore acts both as a target structure that guides extraction and
as a contract that later retrieval and verification can rely on~\cite{hogan2021kg,
pan2024unifying}.

\subsection{Top-Down vs. Bottom-Up Schema Design}

Ontology engineering can proceed top-down, where domain experts define the
classes and relations in advance, or bottom-up, where the schema emerges from the
entities and relations discovered in the data; in practice the two are combined,
because a knowledge graph must balance a principled model with the messiness of
real sources~\cite{hogan2021kg}. The choice affects how tightly extraction is
constrained: a rich predefined schema improves consistency but can miss
unanticipated relations, whereas an open, data-driven schema captures more but
requires later consolidation~\cite{hogan2021kg, ji2021kgsurvey}.

\subsection{Reusing Existing Ontologies}

Because interoperability is a central goal of the Semantic Web, established
vocabularies and ontologies should be reused rather than reinvented wherever
possible, allowing data from different sources to be integrated under shared
semantics~\cite{bernerslee2001semantic}. The RDF and OWL standards exist precisely
to make such vocabularies portable and machine-interpretable across
systems~\cite{w3c2014rdf, w3c2012owl}, and surveys of knowledge graphs document a
large ecosystem of reusable ontologies built on this
foundation~\cite{hogan2021kg}.
